\subsection{OSCILLATION OF A FUNCTION}

Related to the notion of diameter is that of the oscillation of a function \(f\) , defined on an arbitrary set \(\mathbf{X}\) and taking its values in a metric space \(\mathbf{X}'\) ; if \(\mathbf{A}\) is any non-empty subset of \(\mathbf{X}\) , the diameter \(\delta(f(\mathbf{A}))\) is called the oscillation of \(f\) in \(\mathbf{A}\) .

If moreover \(\mathbf{X}\) is a subset of a topological space \(\mathbf{Y}\) , and if \(x \in \overline{\mathbf{X}}\) , the number

\[
\omega (x; f) = \inf \delta (f (\mathbf {V} \cap \mathbf {X}))
\]

(as V runs through the neighbourhood filter of \(x\) in Y) is called the oscillation of \(f\) at \(x \in \overline{\mathbf{X}}\) .

PROPOSITION 4. The oscillation \(\omega(x; f)\) of an arbitrary function \(f\) , defined on a subset X of a topological space Y and taking its values in a metric space \(\mathbf{X}'\) , is upper semi-continuous on \(\overline{\mathbf{X}}\) .

Let \(a\) be any point of \(\overline{\mathbf{X}}\) ; then for each \(k > \omega(a; f)\) there exists an open neighbourhood \(\mathbf{V}\) of \(a\) such that \(\delta(f(\mathbf{V} \cap \mathbf{X})) \leqslant k\) ; for each \(x \in \mathbf{V} \cap \overline{\mathbf{X}}\) , \(\mathbf{V}\) is a neighbourhood of \(x\) and therefore

\[
\omega (x; f) \leqslant \delta (f (\mathbf {V} \cap \mathbf {X})) \leqslant k,
\]

which shows that \(\omega\) is upper semi-continuous at the point a.

In order that \(\omega(x; f) = 0\) at a point \(x \in \overline{\mathbf{X}}\) it is necessary and sufficient that for each \(\varepsilon > 0\) there should exist a neighbourhood \(\mathbf{V}\) of \(x\) such that \(f(\mathbf{V} \cap \mathbf{X})\) is contained in a ball of radius \(\varepsilon\) ; if \(x \in \mathbf{X}\) , this condition expresses the fact that \(f\) is continuous at the point \(x\) (with respect to \(\mathbf{X}\) ); if \(x \in \overline{\mathbf{X}} \cap \complement \mathbf{X}\) , the image under \(f\) of the trace on \(\mathbf{X}\) of the neighbourhood filter of \(x\) in \(\mathbf{Y}\) is a Cauchy filter base on \(\mathbf{X}'\) ; in particular:

PROPOSITION 5. Let \(f\) be a function defined on a subset \(\mathbf{X}\) of a topological space \(\mathbf{Y}\) , taking its values in a complete metric space \(\mathbf{X}'\) . Then \(f\) has a limit relative to \(\mathbf{X}\) at a point \(x \in \overline{\mathbf{X}}\) if and only if the oscillation of \(f\) at \(x\) is zero.

